% !TEX root = Lecture_Koszul.tex
\clearpage
\section{Ext and Tor}\label{sec:ext-tor}

The contravariant functor $\Hom_A(-,N)$ is left exact (Lemma~\ref{lem:Hom-contravariant-left-exact}) but in general not exact (Example~\ref{ex:Hom-contravariant-not-exact}).
The group $\Ext_A^h(M,N)$ measure this failure at homological degree $h$; it is computed by taking homology on the complex  $\Hom_A(P_\bullet,N)$ for a projective resolution $P_\bullet$ of $M$.
We treat $\Ext$ in detail; $\Tor$ is obtained in a similar way from replacing $\Hom_A(-,N)$ by $-\otimes_AN$, and we only record what we need.

\subsection{Ext-groups via projective resolutions}

All modules in this subsection are right $A$-modules over an arbitrary algebra $A$, except where a left module is specified.

Let $p_0:P_0\to M$ and $d_h:P_h\to P_{h-1}$ be a projective resolution $P_\bullet\to M$ (Definition~\ref{def:projective-resolution}), and let $N$ be a right $A$-module.
Postcomposition with the differentials gives the cochain complex $\Hom_A(P_\bullet,N)$, concentrated in degrees $h\geq0$:
\begin{equation}\label{eq:hom-complex}
\Hom_A(P_0,N)\xrightarrow{\delta^0}\Hom_A(P_1,N)\xrightarrow{\delta^1}\Hom_A(P_2,N)\longrightarrow\cdots,
\qquad\delta^h(\varphi)=\varphi\circ d_{h+1}.
\end{equation}
Indeed $\delta^{h+1}\delta^h(\varphi)=\varphi\, d_{h+1}d_{h+2}=0$.
The module $M$ does not appear in this complex.

\begin{definition}\label{def:ext}
For right $A$-modules $M,N$ and a projective resolution $P_\bullet\to M$, put
\[
\Ext_A^h(M,N):=H^h\bigl(\Hom_A(P_\bullet,N)\bigr)\qquad\text{for all }h\geq0.
\]
\end{definition}

This depends only on $M$ and $N$, as the next lemma shows.
It compares a projective resolution of $M$ with an arbitrary resolution of $M'$, which is more than we need now but costs nothing.

\begin{lemma}[Comparison lemma]\label{lem:comparison}
Let $P_\bullet\to M$ be a projective resolution with differentials $d_h$ and augmentation $p_0$.
Let $\cdots\to Q_1\xrightarrow{d'_1}Q_0\xrightarrow{q_0}M'\to0$ be an exact sequence of $A$-modules, and let $f:M\to M'$ be $A$-linear.
\begin{enumerate}
\item There are $A$-linear maps $f_h:P_h\to Q_h$ for $h\geq0$ making the following diagram commute:
\[
\xymatrix@C=26pt@R=26pt{
\cdots\ar[r]&P_2\ar[r]^{d_2}\ar@{-->}[d]^{f_2}&P_1\ar[r]^{d_1}\ar@{-->}[d]^{f_1}&P_0\ar[r]^{p_0}\ar@{-->}[d]^{f_0}&M\ar[r]\ar[d]^{f}&0\\
\cdots\ar[r]&Q_2\ar[r]_{d'_2}&Q_1\ar[r]_{d'_1}&Q_0\ar[r]_{q_0}&M'\ar[r]&0
}
\]
\item If $(f'_h)_{h\geq0}$ is another such family, there are $A$-linear maps $s_h:P_h\to Q_{h+1}$ for $h\geq0$ (the diagonal arrows below) with $f_h-f'_h=d'_{h+1}s_h+s_{h-1}d_h$ for all $h\geq0$, where $s_{-1}:=0$:
\[
\xymatrix@C=34pt@R=34pt{
\cdots\ar[r]&P_2\ar[r]^{d_2}\ar[d]|{f_2-f'_2}&P_1\ar[r]^{d_1}\ar[d]|{f_1-f'_1}\ar[dl]_{s_1}&P_0\ar[d]|{f_0-f'_0}\ar[dl]_{s_0}\\
\cdots\ar[r]&Q_2\ar[r]_{d'_2}&Q_1\ar[r]_{d'_1}&Q_0
}
\]
\end{enumerate}
\end{lemma}

We say that $(f_h)_{h\geq0}$ \defn{lifts} $f$, and that two lifts are \defn{homotopic} when maps $s_h$ as in (2) exist.
For $M=M'$, $Q_\bullet=P_\bullet$ and $f=\id_M$ this recovers the map of resolutions in Theorem~\ref{thm:minimal-resolution}.

\begin{supplementary}
\begin{proof}
(1) Since $P_0$ is projective and $q_0$ is surjective, there is $f_0$ with $q_0f_0=fp_0$.
Suppose $f_0,\ldots,f_{h-1}$ are constructed, $h\geq1$, and consider $f_{h-1}d_h:P_h\to Q_{h-1}$.
For $h=1$, we have $q_0f_0d_1=fp_0d_1=0$, so $f_0d_1$ has image in $\Ker q_0=\Im d'_1$.
For $h\geq2$, $d'_{h-1}f_{h-1}d_h=f_{h-2}d_{h-1}d_h=0$, so $f_{h-1}d_h$ has image in $\Ker d'_{h-1}=\Im d'_h$.
Since $P_h$ is projective and $d'_h:Q_h\to\Im d'_h$ is surjective, there is $f_h:P_h\to Q_h$ with $d'_hf_h=f_{h-1}d_h$.

(2) Put $g_h=f_h-f'_h$, so $q_0g_0=0$ and $d'_hg_h=g_{h-1}d_h$ for $h\geq1$.
The image of $g_0$ lies in $\Ker q_0=\Im d'_1$, so projectivity of $P_0$ gives $s_0$ with $d'_1s_0=g_0$.
Let $h\geq1$ and suppose $s_0,\ldots,s_{h-1}$ are constructed.
Then
\[
d'_h(g_h-s_{h-1}d_h)=g_{h-1}d_h-(g_{h-1}-s_{h-2}d_{h-1})d_h=0,
\]
where $s_{-1}:=0$ and $d_{h-1}d_h=0$.
Thus $g_h-s_{h-1}d_h$ has image in $\Ker d'_h=\Im d'_{h+1}$, and projectivity of $P_h$ gives $s_h$ with $d'_{h+1}s_h=g_h-s_{h-1}d_h$.
\end{proof}
\end{supplementary}

\begin{proposition}\label{prop:ext-well-defined}
Let $P_\bullet\to M$ and $Q_\bullet\to M'$ be projective resolutions, $f:M\to M'$ be $A$-linear, $(f_h)$ a lift of $f$, and $N$ a module.
\begin{enumerate}
\item The map $[\varphi]\mapsto[\varphi f_h]$ is a well-defined linear map $f^*:H^h(\Hom_A(Q_\bullet,N))\to H^h(\Hom_A(P_\bullet,N))$ that does not depend on the choice of lift.
\item If $g:M'\to M''$ is $A$-linear and $R_\bullet\to M''$ is a projective resolution, then $(gf)^*=f^*g^*$; also $(\id_M)^*=\id$.
\end{enumerate}
In particular, up to canonical isomorphism $\Ext_A^h(M,N)$ does not depend on the projective resolution of $M$.
It is contravariant in $M$ through $f^*$.
It is covariant in $N$: an $A$-linear map $\psi:N\to N'$ gives $\psi_*[\varphi]=[\psi\varphi]$.
\end{proposition}

\begin{proof}
(1) Let $\varphi:Q_h\to N$ be a cocycle, so $\varphi d'_{h+1}=0$.
Then $\varphi f_hd_{h+1}=\varphi d'_{h+1}f_{h+1}=0$, so $\varphi f_h$ is a cocycle.
If $\varphi=\psi d'_h$ is a coboundary, then $\varphi f_h=\psi f_{h-1}d_h$ is a coboundary.
For another lift $(f'_h)$ with homotopy $(s_h)$ as in Lemma~\ref{lem:comparison}(2),
\[
\varphi(f_h-f'_h)=\varphi d'_{h+1}s_h+\varphi s_{h-1}d_h=(\varphi s_{h-1})d_h,
\]
a coboundary; for $h=0$ the same computation gives $\varphi d'_1s_0=0$.
(2) If $(f_h)$ lifts $f$ and $(g_h)$ lifts $g$, then $(g_hf_h)$ lifts $gf$, and $(\id_{P_h})$ lifts $\id_M$; use (1).
For two resolutions $P_\bullet,P'_\bullet$ of $M$, the lifts of $\id_M$ in each direction give maps $\id^*$ that are inverse to each other by (2).
Finally $\psi\varphi$ is a cocycle (coboundary) if $\varphi$ is, and $\psi_*$ is compatible with $f^*$.
\end{proof}

\begin{proposition}\label{prop:ext-basic}
For all modules $M,N$:
\begin{enumerate}
\item $\Ext_A^0(M,N)\cong\Hom_A(M,N)$.
\item If $M$ is projective, then $\Ext_A^h(M,N)=0$ for all $h\geq1$.
\end{enumerate}
\end{proposition}

\begin{proof}
(1) Apply Lemma~\ref{lem:Hom-contravariant-left-exact} to the exact sequence $P_1\xrightarrow{d_1}P_0\xrightarrow{p_0}M\to0$.
It shows that $-\circ p_0:\Hom_A(M,N)\to\Hom_A(P_0,N)$ is injective with image $\Ker\delta^0$, the cocycles in degree $0$.
(2) Use the resolution with $P_0=M$ and $P_h=0$ for $h\geq1$.
\end{proof}

Thus $\Ext$ extends $\Hom$ to higher degrees, and it vanishes for projective $M$ as one expects from Proposition~\ref{prop:projective-characterisations}.

\subsection{Graded Ext and minimal resolutions}\label{subsec:graded-ext}

From now on $A$ is as in Subsection~\ref{subsec:graded-resolutions}: locally finite-dimensional non-negatively graded with $S=A_0$ finite-dimensional semisimple, and $S=A/A_+$ is regarded as a graded module in degree $0$.
Write $S=\bigoplus_{y=1}^n\Bbbk e_y$ and $L_y=e_yS$ for the simple modules.
Let $M$ be a bounded below locally finite-dimensional graded right module and $N$ a graded right module.
Let $P_\bullet\to M$ be a graded projective resolution, for instance a minimal one by Theorem~\ref{thm:minimal-resolution}.

\begin{definition}\label{def:graded-ext}
For $h\geq0$ and $j\in\Z$, put
\[
\Ext_A^h(M,N)_j:=H^h\bigl(\Hom_A^{\Z}(P_\bullet,N(j))\bigr).
\]
\end{definition}

Here $\Hom_A^\Z(P_h,N(j))$ is the space of homogeneous maps of degree $j$ (Notation~\ref{not:graded-hom}), and $\delta^h(\varphi)=\varphi\circ d_{h+1}$ as in \eqref{eq:hom-complex}.
The proofs of Lemma~\ref{lem:comparison} and Proposition~\ref{prop:ext-well-defined} apply verbatim with degree-preserving maps, using that graded projective modules lift along degree-preserving surjections and that kernels of graded maps are graded.
Hence $\Ext_A^h(M,N)_j$ is independent of $P_\bullet$ up to canonical isomorphism.
The index $j$ is the \defn{internal degree} or \defn{Adams degree}; the index $h$ is the \defn{homological degree}.

\begin{remark}\label{rem:graded-ext-ungraded}
A graded projective module is projective as an ungraded module, because it is a graded direct summand of a graded free module (Corollary~\ref{cor:graded-projective-characterisations}).
So $P_\bullet\to M$ is also an ordinary projective resolution.
Each $\Hom_A^\Z(P_h,N(j))$ is contained in $\Hom_A(P_h,N)$, and if $P_h$ is finitely generated, then
\[
\Hom_A(P_h,N)=\bigoplus_{j\in\Z}\Hom_A^\Z\bigl(P_h,N(j)\bigr).
\]
Hence, if $P_h$ is finitely generated for all $h$ (for instance if $A$ is right noetherian and $M$ is finitely generated, by Corollary~\ref{cor:noetherian-minimal-resolution}), then the ordinary $\Ext_A^h(M,N)=\bigoplus_{j\in\Z}\Ext_A^h(M,N)_j$.
\end{remark}

\begin{theorem}\label{thm:ext-minimal}
Let $P_\bullet\to M$ be a minimal graded projective resolution and $N$ a graded right $A$-module with $NA_+=0$.
Then every differential of $\Hom_A^\Z(P_\bullet,N(j))$ is zero, and for all $h\geq0$ and $j\in\Z$,
\[
\Ext_A^h(M,N)_j\cong\Hom_A^{\Z}(P_h,N(j))\cong\Hom_S^\Z\bigl(\grtop P_h,N(j)\bigr).
\]
In particular, if $N$ is a right $S$-module placed in degree $0$, then $\Ext_A^h(M,N)_j\cong\Hom_S\bigl((\grtop P_h)_{-j},N\bigr)$.
\end{theorem}

\begin{proof}
Let $\varphi:P_h\to N(j)$ be $A$-linear and degree-preserving.
Then $\varphi(P_hA_+)=\varphi(P_h)A_+\subseteq NA_+=0$, so $\varphi$ factors uniquely through $P_h\to P_h/P_hA_+=\grtop P_h$, and every degree-preserving $S$-linear map $\grtop P_h\to N(j)$ arises this way.
This gives the second isomorphism.
By \eqref{eq:graded-minimal-criterion}, $\Im d_{h+1}\subseteq P_hA_+$, so $\varphi d_{h+1}=0$.
Hence all $\delta^h$ vanish and the cohomology equals the cochain groups.
If $N$ lies in degree $0$, then $N(j)$ lies in degree $-j$, so a degree-preserving map from $\grtop P_h$ only involves $(\grtop P_h)_{-j}$.
\end{proof}

Thus the graded top of $P_h$ is computed by $\Ext$ with values in $S$-modules.
For the simple modules the answer is a multiplicity.

\begin{corollary}\label{cor:ext-simple}
In the setting of Theorem~\ref{thm:ext-minimal}, for $1\leq y\leq n$,
\[
\dim_\Bbbk\Ext_A^h(M,L_y)_j=\bigl[(\grtop P_h)_{-j}:L_y\bigr].
\]
\end{corollary}

\begin{proof}
For a right $S$-module $T$ we have $T\cong\bigoplus_x L_x^{\oplus[T:L_x]}$, and $\Hom_S(L_x,L_y)$ is $\Bbbk$ for $x=y$ and $0$ otherwise.
\end{proof}

\begin{advanced}
\begin{remark}[General $S$]\label{rem:ext-general-S}
Corollary~\ref{cor:ext-simple} used $S=\bigoplus_y\Bbbk e_y$.
For an arbitrary finite-dimensional semisimple $S$, the space $\Hom_S(L_x,L_y)$ is $0$ for $x\neq y$ and the division algebra $\End_S(L_y)$ for $x=y$, so the formula holds with $\dim_\Bbbk$ replaced by the dimension over $\End_S(L_y)$.
\end{remark}
\end{advanced}

The multiplicities $[(\grtop P_h)_i:L_y]$ are the graded Betti numbers $\beta_{h,i,y}(M)$ of Remark~\ref{not:graded-betti}.
Since $S$ is semisimple, $\Hom_S(-,N)$ is exact, so the groups $\Ext_A^h(M,N)_j$ with $h\geq1$ record the resolution over $A$, not any further structure of $S$.

\begin{example}\label{ex:ext-A2}
Let $A=\Bbbk A_2$.
By Example~\ref{ex:minimal-A2}, the minimal resolution of $L_1$ is $0\to e_{22}A(-1)\to e_{11}A\to L_1\to0$, so $\grtop P_0=L_1$ in degree $0$ and $\grtop P_1=L_2$ in degree $1$.
By Corollary~\ref{cor:ext-simple}, the only non-zero groups $\Ext_A^h(L_1,L_y)_j$ are
\[
\Ext_A^0(L_1,L_1)_0\cong\Bbbk,\qquad\Ext_A^1(L_1,L_2)_{-1}\cong\Bbbk.
\]
The class in $\Ext_A^1(L_1,L_2)_{-1}$ reflects the exact sequence $0\to L_2\to e_{11}A\to L_1\to0$, which does not split since $e_{11}A$ is not semisimple.
In general, $\Ext_A^1(M,N)$ classifies such extensions of $M$ by $N$; we do not use this.
\end{example}

\begin{example}\label{ex:ext-polynomial-exterior}
For $A=\Bbbk[x]$ with $\deg x=1$, Example~\ref{ex:minimal-polynomial} gives $\grtop P_0=\Bbbk$ in degree $0$ and $\grtop P_1=\Bbbk$ in degree $1$, so
\[
\Ext_A^0(\Bbbk,\Bbbk)_0\cong\Bbbk,\qquad\Ext_A^1(\Bbbk,\Bbbk)_{-1}\cong\Bbbk,
\]
and all other groups vanish.
For $E=\Bbbk[\epsilon]/(\epsilon^2)$ with $\deg\epsilon=1$, Example~\ref{ex:minimal-exterior} gives $\grtop P_h=\Bbbk$ in degree $h$ for all $h\geq0$, so $\Ext_E^h(\Bbbk,\Bbbk)_j\cong\Bbbk$ if $j=-h$ and $0$ otherwise.
\end{example}

In all three examples the non-zero groups $\Ext_A^h(S,S)_j$ lie on the line $j=-h$, because the generators of $P_h$ have degree $h$; we return to this in Section~\ref{sec:koszul-complex}.

\subsection{Tor}

Let $M$ be a right $A$-module with a projective resolution $P_\bullet\to M$, and $N$ a left $A$-module.
The maps $d_h\otimes\id_N$ make $P_\bullet\otimes_AN$ a chain complex of vector spaces (Definition~\ref{def:chain-complex}).

\begin{definition}\label{def:tor}
$\Tor_h^A(M,N):=H_h(P_\bullet\otimes_AN)$ for all $h\geq0$.
\end{definition}

As for $\Ext$, this is independent of $P_\bullet$ up to canonical isomorphism and functorial in $M$ and $N$, with $\Tor_0^A(M,N)=M\otimes_AN$; for graded $M,N$ and a graded projective resolution it is graded.

\begin{supplementary}
The proofs of Lemma~\ref{lem:comparison} and Proposition~\ref{prop:ext-well-defined} apply with $\Hom_A(-,N)$ replaced by $-\otimes_AN$, which sends lifts and homotopies to lifts and homotopies.
Right exactness of $-\otimes_AN$ gives $\Tor_0^A(M,N)=M\otimes_AN$.
The grading is $\Tor_h^A(M,N)_i=H_h((P_\bullet\otimes_AN)_i)$ for the total grading of Definition~\ref{def:total-tensor-grading}.
\end{supplementary}

We only need $N=S=A/A_+$, a graded left $A$-module in degree $0$.

\begin{proposition}\label{prop:tor-top}
Let $P_\bullet\to M$ be a minimal graded projective resolution.
Then $\Tor_h^A(M,S)\cong\grtop P_h$ as graded right $S$-modules, for all $h\geq0$.
\end{proposition}

\begin{supplementary}
\begin{proof}
We have $P_h\otimes_AS=P_h/P_hA_+=\grtop P_h$, and the differential of $P_\bullet\otimes_AS$ is induced by $d_h$.
It is zero because $\Im d_h\subseteq P_{h-1}A_+$.
\end{proof}
\end{supplementary}

For a graded left $A$-module $N$, the graded linear dual $\bbD(N)$ is a right $A$-module via $(fa)(n)=f(an)$, and $\bbD(N)_i=D(N_{-i})$.

\begin{proposition}\label{prop:ext-tor-duality}
For graded $M,N$ as above, with $N$ a graded left $A$-module, and all $h\geq0$ and $j\in\Z$,
\[
\Ext_A^h\bigl(M,\bbD(N)\bigr)_j\cong D\bigl(\Tor_h^A(M,N)_{-j}\bigr).
\]
\end{proposition}

\begin{supplementary}
\begin{proof}
A degree-preserving $A$-linear map $\varphi:P\to\bbD(N)(j)$ is the same as a bilinear form $b(p,n)=\varphi(p)(n)$ on $P\times N$ with $b(pa,n)=b(p,an)$ that vanishes on $P_i\times N_k$ unless $i+k=-j$.
This is a linear form on $(P\otimes_AN)_{-j}$, so $\Hom_A^\Z(P,\bbD(N)(j))\cong D((P\otimes_AN)_{-j})$, naturally in $P$.
Applying this to $P_\bullet$ identifies the cochain complex $\Hom_A^\Z(P_\bullet,\bbD(N)(j))$ with $D((P_\bullet\otimes_AN)_{-j})$.
Since $D$ is exact, it commutes with taking (co)homology.
\end{proof}
\end{supplementary}

For $N=S$ we have $\bbD(S)\cong S$ as graded right $A$-modules, since the dual basis $(e_x^*)_x$ of $D(S)$ satisfies $e_x^*e_y=\delta_{xy}e_x^*$ and both modules are killed by $A_+$.
Hence, for every bounded below locally finite-dimensional $M$,
\[
\Ext_A^h(M,S)_j\cong D\bigl((\grtop P_h)_{-j}\bigr),
\]
by Propositions~\ref{prop:tor-top} and~\ref{prop:ext-tor-duality}; this agrees with Theorem~\ref{thm:ext-minimal}.
So $\Ext_A^\bullet(M,S)$ and $\Tor_\bullet^A(M,S)$ carry the same information, and we work with $\Ext$.

\begin{advanced}
\begin{remark}[General $S$]\label{rem:tor-general-S}
The identification $\bbD(S)\cong S$ above uses $S=\bigoplus_x\Bbbk e_x$.
For arbitrary semisimple $S$, the module $\bbD(S)$ is isomorphic to $S$ only after choosing a symmetrising form; without it one can use $\Ext_A^h(M,\bbD(S))$ in place of $\Ext_A^h(M,S)$ throughout.
\end{remark}
\end{advanced}

\subsection{The beginning of the resolution of $S$}

We determine the first terms of the minimal resolution of $S$ for algebras generated in degree 1 and for quadratic algebras.
Here $A$ is as in Subsection~\ref{subsec:graded-ext}, and generation in degree $1$ over $S$ is as in \eqref{eq:multiplication-maps}.

\begin{proposition}\label{prop:resolution-first-terms}
Let $P_\bullet\to S$ be the minimal graded projective resolution of $S=A/A_+$.
\begin{enumerate}
\item $P_0=A$, with $p_0:A\to S$ the projection.
\item $\grtop P_1\cong A_+/A_+^2$.
If $A$ is generated in degree $1$, then $\grtop P_1\cong A_1$, concentrated in degree $1$, and $P_1\cong(A_1\otimes_SA)(-1)$ with $d_1(v\otimes a)=va$.
\item If $A$ is generated in degree $1$, put $V=A_1$ and let $I=\Ker(T_S(V)\onto A)$ as in Corollary~\ref{cor:tensor-presentation}.
Then
\[
(\grtop P_2)_i\cong I_i/\bigl(V\otimes_SI_{i-1}+I_{i-1}\otimes_SV\bigr)\quad\text{for all }i.
\]
In particular, if $A=T_S(V)/(R)$ is quadratic with $V=A_1$ and $R\subseteq V\otimes_SV$, then $\grtop P_2\cong R$, concentrated in degree $2$, and $P_2\cong(R\otimes_SA)(-2)$.
\end{enumerate}
In particular $\Tor_h^A(S,S)\cong S,\ A_1,\ R$ for $h=0,1,2$, concentrated in degrees $h$, in the quadratic case.
\end{proposition}

\begin{supplementary}
\begin{proof}
(1) The module $A$ is projective and $\Ker p_0=A_+=AA_+$, so $\Ker p_0\subseteq P_0A_+$ and $p_0$ is the first step of a minimal resolution.

(2) Put $K_1=\Ker p_0=A_+$.
By Lemma~\ref{lem:minimal-tops}, $\grtop P_1\cong\grtop K_1=A_+/A_+^2$.
If $A$ is generated in degree $1$, then $A_+^2=\bigoplus_{i\geq2}A_i$, so $\grtop K_1=A_1$ in degree $1$.
By \eqref{eq:graded-cover-terms}, $P_1\cong(A_1\otimes_SA)(-1)$, and the surjection $P_1\to K_1$ is $v\otimes a\mapsto va$.

(3) Put $T=T_S(V)$, so $A_i=T_i/I_i$ and $I_1=0$.
Since every $S$-module is projective, $V\otimes_S-$ is exact, so
\[
(P_1)_i=V\otimes_SA_{i-1}=T_i/(V\otimes_SI_{i-1}),
\]
and $d_1:(P_1)_i\to A_i=T_i/I_i$ is induced by the identity on $T_i$.
Hence $K:=\Ker d_1$ has $K_i=I_i/(V\otimes_SI_{i-1})$.
Since $A_j=A_1^j$, we have $(KA_+)_i=K_{i-1}A_1$, which is the image of $I_{i-1}\otimes_SV$.
Therefore $(K/KA_+)_i=I_i/(V\otimes_SI_{i-1}+I_{i-1}\otimes_SV)$, and $\grtop P_2\cong\grtop K$ by Lemma~\ref{lem:minimal-tops}.

If $I=(R)$, then $I_i=\sum_{p+q=i-2}T_pRT_q$.
For $i\geq3$ each summand with $p\geq1$ lies in $V\otimes_SI_{i-1}$, and each with $p=0$ lies in $I_{i-1}\otimes_SV$.
So $I_i=V\otimes_SI_{i-1}+I_{i-1}\otimes_SV$ for $i\geq3$, while $I_1=0$ and $I_2=R$.
Thus $\grtop P_2=R$ in degree $2$, and \eqref{eq:graded-cover-terms} gives $P_2\cong(R\otimes_SA)(-2)$.
The last sentence follows from Proposition~\ref{prop:tor-top}.
\end{proof}
\end{supplementary}

Dually, by Theorem~\ref{thm:ext-minimal} the space $\Ext_A^h(S,S)_{-h}$ has dimension $\dim_\Bbbk\grtop P_h$, which is $n$, $\dim_\Bbbk A_1$, $\dim_\Bbbk R$ for $h=0,1,2$ when $A$ is quadratic.
These are the dimensions of $A^!_0=S$, $A^!_1=D(A_1)$ and $A^!_2\cong D(R)$ from Section~\ref{sec:quadratic-dual}, so $\Ext_A^\bullet(S,S)$ and $A^!$ already agree in dimension up to degree $2$; making this precise requires a product on $\Ext$.

\subsection{Exercises}

\begin{exercise}
\begin{enumerate}
\item Prove that $\Ext_A^h(P,N)=0$ for $h\geq1$ and all $N$ if $P$ is projective.
\item Conversely, let $p_0:P_0\to M$ be a surjection from a projective module with $K=\Ker p_0$, and suppose $\Ext_A^1(M,K)=0$.
Prove that $M$ is projective.
\end{enumerate}
\emph{Hint for (2).} Take $P_1\to K$ onto $K$ and consider it as a cocycle in $\Hom_A(P_1,K)$.
\end{exercise}

\begin{hidden}
\emph{Solution.}
(1) $P_0=P$ and $P_h=0$ for $h\geq1$ is a projective resolution of $P$, so $\Hom_A(P_\bullet,N)$ is concentrated in degree $0$.

(2) Take a surjection $\pi:P_1\to K$ from a projective module, and extend $P_1\xrightarrow{\iota\pi}P_0$ to a resolution $P_\bullet\to M$, where $\iota:K\hookrightarrow P_0$.
Since $\Ker(\iota\pi)=\Ker\pi=\Im d_2$, we have $\pi d_2=0$, so $\pi$ is a cocycle in $\Hom_A(P_1,K)$.
As $\Ext_A^1(M,K)=0$, there is $\sigma\in\Hom_A(P_0,K)$ with $\pi=\sigma d_1=\sigma\iota\pi$.
Since $\pi$ is surjective, $\sigma\iota=\id_K$.
So $K$ is a direct summand of $P_0$, with complement $\cong P_0/K\cong M$, and $M$ is projective.
\end{hidden}


\begin{exercise}
Let $P_\bullet\to M$ be a projective resolution and $K_1=\Ker p_0$.
Prove that
\[
\Ext_A^1(M,N)\cong\operatorname{coker}\bigl(\Hom_A(P_0,N)\to\Hom_A(K_1,N)\bigr),
\]
where the map is restriction, and that $\Ext_A^h(M,N)\cong\Ext_A^{h-1}(K_1,N)$ for $h\geq2$.
\end{exercise}

\begin{hidden}
\emph{Solution.}
The shifted complex $\cdots\to P_3\to P_2\to P_1\xrightarrow{d_1}K_1\to0$ is a projective resolution of $K_1$, so $\Ext_A^{h-1}(K_1,N)$ is the cohomology of $\Hom_A(P_{\bullet+1},N)$ in degree $h-1$.

For $h\geq2$, the cohomology of $\Hom_A(P_\bullet,N)$ in degree $h$ uses $\Hom_A(P_{h-1},N)\to\Hom_A(P_h,N)\to\Hom_A(P_{h+1},N)$, which is the same piece of the complex, so $\Ext_A^h(M,N)\cong\Ext_A^{h-1}(K_1,N)$.

For $h=1$: $\Ker\delta^1$ consists of maps $P_1\to N$ vanishing on $\Im d_2=\Ker d_1$, that is, maps factoring uniquely through $d_1:P_1\twoheadrightarrow K_1$.
So $\Ker\delta^1\cong\Hom_A(K_1,N)$, and under this identification $\Im\delta^0=\{\varphi d_1:\varphi\in\Hom_A(P_0,N)\}$ is the image of the restriction map $\Hom_A(P_0,N)\to\Hom_A(K_1,N)$, $\varphi\mapsto\varphi|_{K_1}$.
Hence $\Ext_A^1(M,N)$ is the cokernel.
\end{hidden}


\begin{exercise}
Let $A=\Bbbk A_3/(e_{13})$ with the grading by path length.
Find the minimal graded projective resolution of $L_1$, and compute $\dim_\Bbbk\Ext_A^h(L_1,L_y)_j$ for all $h,y,j$.
\end{exercise}

\begin{hidden}
\emph{Solution.}
Here $e_{11}A=\langle e_{11},e_{12}\rangle$, $e_{22}A=\langle e_{22},e_{23}\rangle$, $e_{33}A=\langle e_{33}\rangle$.
The kernel of $e_{11}A\to L_1$ is $\Bbbk e_{12}\cong L_2(-1)$, since $e_{12}e_{23}=e_{13}=0$.
The map $e_{22}A(-1)\to\Bbbk e_{12}$, $e_{22}\mapsto e_{12}$, is surjective with kernel $\Bbbk e_{23}\subseteq e_{22}A_+$ in degree $2$, and $\Bbbk e_{23}\cong e_{33}A(-2)$ is projective.
So the minimal resolution is
\[
0\to e_{33}A(-2)\to e_{22}A(-1)\to e_{11}A\to L_1\to0,
\]
with $\grtop P_h=L_{h+1}$ in degree $h$ for $h=0,1,2$.
By Corollary~\ref{cor:ext-simple}, $\dim\Ext_A^h(L_1,L_y)_j=[(\grtop P_h)_{-j}:L_y]$ equals $1$ for $(h,y,j)=(0,1,0),(1,2,-1),(2,3,-2)$ and $0$ otherwise.
\end{hidden}


\begin{exercise}
Let $A=\Bbbk[x]$ with $\deg x=1$ and $M=A/(x^m)$ for some $m\geq1$.
Find a minimal graded resolution of $M$, and compute $\Tor_h^A(M,\Bbbk)$ and $\Ext_A^h(M,\Bbbk)_j$ for all $h,j$.
\end{exercise}

\begin{hidden}
\emph{Solution.}
Multiplication by $x^m$ is injective on $A$, so $0\to A(-m)\xrightarrow{x^m}A\to M\to0$ is a graded projective resolution, and it is minimal since $x^m\in A_+$ for $m\geq1$.
So $\grtop P_0=\Bbbk$ in degree $0$ and $\grtop P_1=\Bbbk$ in degree $m$, and $\grtop P_h=0$ for $h\geq2$.
By Proposition~\ref{prop:tor-top}, $\Tor_h^A(M,\Bbbk)\cong\grtop P_h$, which is $\Bbbk$ in degree $0$ for $h=0$, $\Bbbk$ in degree $m$ for $h=1$, and $0$ otherwise.
By Theorem~\ref{thm:ext-minimal}, $\Ext_A^h(M,\Bbbk)_j\cong\Hom_\Bbbk((\grtop P_h)_{-j},\Bbbk)$ is $\Bbbk$ for $(h,j)=(0,0)$ and $(1,-m)$, and $0$ otherwise.
\end{hidden}


\begin{exercise}
Let $A$ be as in Subsection~\ref{subsec:graded-ext}.
\begin{enumerate}
\item Prove that $\Tor_1^A(S,S)\cong A_+/A_+^2$ as graded right $S$-modules.
\item Prove that $A$ is generated in degree $1$ over $S$ if and only if $\Tor_1^A(S,S)$ is concentrated in degree $1$.
\end{enumerate}
\end{exercise}

\begin{hidden}
\emph{Solution.}
(1) Let $P_\bullet\to S$ be the minimal resolution and $K_1=\Ker p_0=A_+$.
By Proposition~\ref{prop:tor-top}, $\Tor_1^A(S,S)\cong\grtop P_1$, and $\grtop P_1\cong\grtop K_1=A_+/A_+^2$ by Lemma~\ref{lem:minimal-tops} applied to $d_1:P_1\to K_1$.

(2) By (1), the condition says $A_+/A_+^2$ is concentrated in degree $1$, that is, $A_i=(A_+^2)_i=\sum_{0<p<i}A_pA_{i-p}$ for $i\geq2$.
If $A$ is generated in degree $1$, then $A_i=A_1A_{i-1}\subseteq A_+^2$ for $i\geq2$.
Conversely, if $A_i=\sum_{0<p<i}A_pA_{i-p}$ for all $i\geq2$, induction on $i$ gives $A_p=A_1^p$ for all $p<i$, hence $A_i=A_1^i$, so every $\mu_i$ is surjective.
\end{hidden}

